神经网络不是``很多矩阵乘法''四个字能够解释的对象。矩阵乘法规定坐标怎样混合，非线性决定能表示什么函数，损失函数定义什么算好，反向传播计算参数改变会怎样影响损失，优化器才真正更新参数。任何一环含混，训练曲线都可能看似正常而模型学错目标。

本单元按这条因果链组织：

\begin{enumerate}
\def\labelenumi{\arabic{enumi}.}
\tightlist
\item
  神经网络为何需要非线性先证明没有激活函数时再深也只是一层线性映射，再讨论 ReLU、sigmoid 等怎样同时改变表示和梯度。
\item
  损失函数决定网络在学习什么从条件均值、概率分布和业务决策出发选择目标，避免把训练 surrogate、评估指标和真实成本混成一个数。
\item
  一次完整的前向与反向传播用标量例子和 batch 矩阵公式走完链式法则，说明自动微分究竟缓存和传播什么。
\item
  训练动力学与数值稳定把梯度送进 mini-batch 优化，连接学习率、初始化、梯度流、稳定概率计算和训练诊断。
\item
  通用逼近定理说清了什么又没说什么接续第一篇的表示追问：先陈述一层非线性网络在紧集上稠密并用台阶直觉解释为什么，再逐条拆穿定理没说的宽度、可优化性与泛化，最后说明深度换宽度的效率来自何处。
\end{enumerate}

本单元的目的不是覆盖所有网络架构，而是建立一套不会随框架 API 改名而失效的分析方法：写清 shape，追踪函数复合，识别损失的统计含义，区分数学条件与浮点实现，分清表示、优化与泛化分别由什么保证，再用小规模可计算例子验证。
